% 第一章 绪论
\thesischapterpage{1\quad 绪论}
\pagenumbering{arabic}

\titleformat{\section}{\heiti\zihao{3}\centering}{\thesection}{0.5em}{}[]
\section{绪论}

\subsection{研究背景与意义}
\thesisbodyfont
随着机械手系统逐步进入家庭服务、柔性装配和开放环境操作场景，手部操作研究的关注点也在变化：以前更多关心“能不能抓住”，现在则更看重“抓住以后还能不能继续处理”。相比平行夹爪，多指灵巧手具有更丰富的运动自由度和接触构型，更适合完成姿态调整、工具使用和精细装配等任务\cite{bicchi2000,okamura2000,openai2018,xia2022}。在这样的背景下，手内操作逐渐成了灵巧操作研究中绕不开的问题之一。

手内旋转是手内操作里最典型、也最能体现连续接触控制难度的基础技能之一。它要求机器手在不重新抓取物体的前提下，仅靠手指之间的协调接触持续调整物体姿态。和一次性的抓取或搬运不同，这个过程里同时包含了高维控制、接触切换、摩擦不确定性和姿态稳定性等问题，因此对策略稳定性、感知质量和控制精度都提出了更高要求。对实际应用来说，手内旋转可以直接服务于工件对位、工具使用和姿态调整等场景，工程价值比较明确。

传统解析模型在这个问题上并不轻松。手内旋转过程中常常伴随复杂的非线性接触、摩擦不确定性以及频繁的接触模式切换，精确建模的代价很高。只要对象质量分布、摩擦特性或外部扰动发生变化，基于固定模型设计的控制器就可能需要重新标定。阻抗控制等经典方法确实为接触稳定性分析提供了重要基础\cite{hogan1985}，但在高维多指接触和对象参数变化明显的条件下，仍然离不开大量建模和调参工作。近年来，深度强化学习提供了另一条路：通过与环境的大规模交互，策略可以在高保真仿真中逐步学得连续动作和长期稳定性，再借助域随机化等方法提升对现实差异的适应能力\cite{tobin2017,peng2018,openai2018,openai2019,isaacgym2021}。

尽管如此，现有成果仍主要集中在高成本平台、特定对象或附加感知条件下的手内操作研究。对于以 LEAP Hand 为代表的低成本灵巧手平台，如何围绕连续手内旋转任务构建一整套兼顾训练效率、部署可行性和后续扩展潜力的方案，仍然是一个很现实的问题\cite{leaphand2023}。和 Shadow Hand 这类高成本平台相比，LEAP Hand 结构紧凑、成本较低、容易复现实验，但也面临电机输出余量有限、关节控制精度受限、缺少高精度触觉传感以及真实接触参数难以准确标定等局限。这些硬件约束使策略在部署阶段难以依赖丰富外部感知或精确动力学模型，而更需要从关节状态、关节目标及其时间变化中恢复隐含的接触和对象信息。因此，本文采用以本体感觉历史为核心的设计思路，并在深度强化学习框架下完成连续旋转策略训练与真实部署研究。

因此，本文的价值主要体现在两个具体层面：一是验证低成本 LEAP Hand 是否能够通过仿真强化学习获得可部署的连续旋转能力；二是把任务建模、两阶段学习和 Sim2Real 部署中的关键接口整理成一条可复现实验链路，为后续扩展到更复杂对象提供基础。

\subsection{国内外研究现状}
\thesisbodyfont
在灵巧手手内操作研究中，强化学习与 Sim2Real 结合是一条重要路线。OpenAI 在 Shadow Hand 平台上完成了基于视觉的手内重定向策略学习，展示了仿真训练结合域随机化在复杂多指系统中的迁移潜力\cite{openai2018}。随后，OpenAI 又展示了单手解魔方任务，将长期规划、精细接触操作与零样本仿真到现实迁移结合起来，使高复杂度灵巧操作受到更多关注\cite{openai2019}。

从训练方法和实验平台的发展来看，PPO 由于实现简洁、训练稳定，在机器人强化学习任务中被广泛采用\cite{ppo2017}。与此同时，Isaac Gym 这类 GPU 并行仿真平台显著提高了样本采集效率，使得在单机环境下训练大规模灵巧手策略成为可能\cite{isaacgym2021}。强化学习算法与高并行仿真平台的结合，推动了灵巧操作研究从单一演示逐步走向系统化训练与评估。

具体到灵巧手学习控制，Rajeswaran 等将示范数据与深度强化学习结合，用于提升高维灵巧操作任务的学习效果\cite{rajeswaran2018}；Zhu 等进一步展示了在低成本硬件上通过强化学习完成灵巧操作的可能性\cite{zhu2019}；Van Hoof 等则较早探索了触觉特征在机器人手内操作学习中的作用\cite{vanhoof2015}。这些工作分别从示范利用、低成本平台和触觉感知三个角度推进了学习式灵巧操作的发展。近年来，手内旋转任务也出现了若干代表性研究。Qi 等提出 HORA，通过在圆柱体对象上训练基础策略并利用本体感觉历史进行在线适配，实现了对多种真实物体的直接旋转部署，表明基于快速自适应的两阶段方法在手内旋转任务中具有较高实用价值\cite{hora2022}。在此基础上，Qi 等又提出融合视觉和触觉信息的 RotateIt，将研究对象扩展到多轴旋转与多模态感知条件下的手内旋转，进一步提升了对复杂对象属性的识别与控制能力\cite{rotateit2023}。Yang 等提出 AnyRotate，利用密集触觉信息实现了重力方向变化条件下的多轴手内旋转，并展示了面向未见对象的零样本迁移能力\cite{anyrotate2024}。代表性工作的横向对比如表 \ref{tab:chapter1_related_work_compare} 所示。

\begin{table}[htbp]
    \centering
    \zihao{5}
    \caption{手内旋转相关代表性工作的横向对比}
    \label{tab:chapter1_related_work_compare}
    \setlength{\tabcolsep}{3pt}
    \renewcommand{\arraystretch}{1.18}
    \begin{tabular}{
        >{\raggedright\arraybackslash}p{0.17\textwidth}
        >{\raggedright\arraybackslash}p{0.17\textwidth}
        >{\raggedright\arraybackslash}p{0.19\textwidth}
        >{\raggedright\arraybackslash}p{0.16\textwidth}
        >{\raggedright\arraybackslash}p{0.19\textwidth}
    }
        \toprule
        工作 & 平台/对象 & 感知模态 & 任务特点 & 迁移或泛化方式 \\
        \midrule
        OpenAI手内操作\cite{openai2018,openai2019} & Shadow Hand，高成本平台 & 视觉、本体感觉等 & 姿态重定向、魔方操作 & 大规模域随机化，真实部署 \\
        HORA\cite{hora2022} & 灵巧手，圆柱体基础对象 & 本体历史、训练期特权信息 & 单轴连续旋转 & 两阶段快速运动适应 \\
        RotateIt\cite{rotateit2023} & 多类旋转对象 & 视觉、触觉、本体感觉 & 多轴手内旋转 & 多模态估计对象属性 \\
        AnyRotate\cite{anyrotate2024} & 多对象，重力方向变化 & 密集触觉、本体感觉 & 多轴旋转，未见对象 & 触觉驱动零样本迁移 \\
        \bottomrule
    \end{tabular}
\end{table}

表 \ref{tab:chapter1_related_work_compare} 中的对比显示，HORA 已经证明两阶段快速适应路线在手内旋转任务中的价值，RotateIt 和 AnyRotate 则把研究推进到多模态感知、多轴旋转和更复杂对象泛化上。本文并不以“全面超过这些工作在任意物体泛化上的能力”为目标，而是聚焦一个更具体的缺口：在低成本 LEAP Hand、无视觉/无触觉反馈、仅依赖本体感觉历史的条件下，建立一条可复现的圆柱体连续旋转训练、评估与实机闭环链路，并分析该链路在圆柱体参数分布外条件下的主要退化因素。

国内在灵巧操作领域的研究起步相对较晚，但近年来在感知、抓取与泛化学习方面进展明显。Xia 等对灵巧操作中的感知问题进行了系统综述，指出视觉、触觉和本体感觉的协同是提升灵巧操作能力的重要方向\cite{xia2022}。在学习方法方面，Yuan 等提出跨构型灵巧抓取强化学习框架，探索了统一动作表示在不同灵巧手平台之间的迁移能力，反映出国内研究已经开始从单一平台的策略学习转向更强调泛化与统一表示的方向\cite{yuan2025}。不过，与灵巧抓取和多模态感知研究相比，直接面向连续手内旋转并强调完整 Sim2Real 部署链路的工作仍然不多。

已有研究已经把手内旋转推进到多模态感知、多轴旋转和未见对象迁移等方向，但这些成果往往依赖更昂贵的平台、更强的传感配置，或更复杂的对象状态估计。低成本灵巧手上的连续旋转仍有一个实际缺口：在视觉和触觉都不作为策略输入的情况下，系统能否仅凭本体感觉历史完成稳定闭环，并在圆柱体参数变化时保持一定鲁棒性。本文聚焦这一缺口，研究范围主要限定在圆柱体及少量近圆柱对象，不与多模态方法在任意对象泛化能力上作直接比较。

\subsection{本文研究内容与主要贡献}
\thesisbodyfont
基于上述研究现状，本文将研究重点收敛到 LEAP Hand 平台上的圆柱体连续手内旋转任务，目标是在 Isaac Lab 仿真环境中建立可训练、可评估、可部署的连续旋转系统，并在真实平台上完成验证。围绕这一目标，本文主要回答三个仍未被现有文献直接覆盖的问题。第一，HORA 类两阶段适应方法在低成本 LEAP Hand、无视觉和无触觉反馈、真实执行器余量有限的条件下，是否仍然能够形成可执行的连续旋转闭环。第二，训练期能够获得的对象位置、尺度、质量、摩擦和质心偏移等属性，应该直接作为策略主干的输入，还是更适合压缩为面向控制的低维技能先验，由本体感觉历史来估计。第三，当第二阶段部署策略在圆柱体物理参数分布外条件下出现退化时，主要风险来自哪些对象属性，教师——学生动作一致性约束和质心偏移覆盖能否缓解这类退化。

围绕这三个问题，本文按“任务建模——策略学习——部署验证”的逻辑展开。首先，结合 LEAP Hand 的运动特性和真实控制接口，建立圆柱体连续旋转任务的状态观测、动作表示、奖励函数、终止条件和评测指标，使仿真训练、统一评估与实机部署具有一致定义。其次，构建技能先验驱动的本体自适应旋转框架：第一阶段利用训练期可获得的对象物理信息学习旋转技能先验和教师策略，第二阶段在冻结动作主干的基础上，通过本体感觉历史估计同一技能先验，从而让部署策略在无法直接获得对象质量、摩擦和质心等参数的条件下仍能闭环执行。最后，针对第二阶段在分布外退化的问题，本文引入教师——部署动作一致性约束，并扩展质心偏移等风险参数的训练覆盖，令部署侧在保留低维技能先验的同时尽量贴近教师侧动作语义。

在仿真验证层面，本文完成 ID/OOD 圆柱体参数范围下的量化评测，并通过消融实验分析特权信息类型、历史长度、历史编码器结构与动作一致性损失的作用。实机部分将同一增强策略部署到真实 LEAP Hand，在标准圆柱体、尺寸变化圆柱体和若干近圆柱真实物体上进行闭环测试，用真实执行结果检查策略链路的可执行性与边界。相应地，本文的主要贡献包括：面向 LEAP Hand 连续手内旋转建立与真实部署接口一致的任务建模、奖励设计和统一评测流程；构建技能先验驱动的两阶段本体自适应框架，把训练期特权信息学习到的旋转能力转化为仅依赖本体感觉历史的可部署策略；针对第二阶段分布外退化，进一步引入动作一致性约束和风险参数覆盖，并在 ID/OOD 仿真评测及真实 LEAP Hand 闭环部署中验证其在当前任务范围内的有效性。

\subsection{论文结构}
\thesisbodyfont
全文共分为五章。第一章介绍研究背景、相关工作、本文的研究内容与论文结构。第二章围绕连续手内旋转任务进行建模，说明任务定义、状态与动作表示、奖励设计以及训练目标。第三章介绍算法设计，重点说明技能先验驱动的本体自适应旋转框架及其部署一致性设计。第四章给出实验配置、仿真评测、实机部署结果及综合分析。第五章总结全文工作，并对后续研究方向进行展望。

\subsection{本章小结}
\thesisbodyfont
本章介绍了连续手内旋转任务的研究背景、国内外研究现状、本文研究内容和论文结构。相关工作已经将手内旋转研究推进到泛化能力和真实部署阶段，但低成本灵巧手平台上的完整训练与部署链路仍不充分。基于这一判断，后续章节将围绕 LEAP Hand 连续手内旋转任务展开建模、方法设计与实验验证。
